\section{Introduction}

Language-model agents increasingly act on the web on a user's behalf: they search, call APIs and
assemble results. Paying for those calls has required accounts, cards and invoices that an
autonomous agent cannot hold. The x402 protocol~\cite{x402} removes that barrier. A server answers a
request with HTTP status 402 and a machine-readable price, the client returns a signed stablecoin
authorization~\cite{eip3009}, and settlement on a blockchain takes seconds.

An agent that buys this way must do two things. It must read what the seller sends (price, payee and
a free-text description) to decide whether a purchase serves its goal, and it must then produce a
payment. In the common design both happen in one process, and that process holds the signing
key. This is where \emph{indirect prompt injection}~\cite{greshake2023} does the most damage: text a
seller places in its own description can steer the agent, and a steered agent with a payment key
moves money that cannot be recovered. A 2026 code review, summarised by the Cloud Security
Alliance~\cite{csa2026}, reported that default x402 clients applied no spending limit until fixes in
August 2026, and the per-payment cap those fixes added can be switched off by a flag.

Separating the model's reasoning from spending authority is not new. CaMeL~\cite{camel} checks
capability policies outside the model that reads untrusted data, and PACE~\cite{pace} has a verifier
sign each transaction a smart account will execute. We take that separation as given and ask a
narrower question: when the agent's process is fully compromised, what spending authority does it keep
over the life of a goal, and what does it learn about the user's limits? Answering it for x402 turned
up three problems that a design for this setting has to solve, and they organise this paper.

The first is that a signature outlives the approval behind it. An EIP-3009 transfer signed for an
approved request can be collected whatever later happens to the approval, so the signer must be right
about the chain: a lying RPC endpoint, or an approval on a block that is later dropped, becomes money.
We remove that dependence by signing only payments that x402's auth-capture escrow can collect, with
the policy contract as the escrow's operator, so the contract's own checks run on the canonical chain
when money moves (Section~\ref{sec:escrow}). The second is setup. The user funds an address that
reaches the console through the agent's process, so a compromised agent can name an address it holds.
The payer is now derived for the user's wallet inside the signer, and a goal opens only with the
signer's statement saying so (Section~\ref{sec:enrollment}). The third is leakage. Refunding an approval
that was never paid is necessary for accounting and makes the policy free to question: a request that
consumes nothing still has its decision published. A private policy then needs a separate budget of
questions, and even with one, the private limit protects less than its encryption suggests
(Section~\ref{sec:leakeval}).

\paragraph{Contributions.}
\begin{enumerate}
  \item \emph{Escrow-gated authorization with authenticated setup, and a lifetime spending bound.}
        Under an adversary with arbitrary code execution in the agent, on a goal whose setup and
        signer are honest, vendor payments over the goal's life total at most
        $\min(\text{budget}, \text{cap} \times \text{calls})$, to allowlisted payees only, whatever the
        signer's RPC reports. We state the bound over an explicit transition model, prove the
        invariants it rests on for every transition, and test them by stateful fuzzing against the
        deployed escrow's code, by mutation, and by a live adversary campaign
        (Sections~\ref{sec:design}--\ref{sec:security} and~\ref{sec:eval}).
  \item \emph{Refund-induced leakage.} Reclaim, which restores spending capacity to unpaid
        approvals, let an adversary binary-search the encrypted cap of our own deployed contract without
        permanently consuming budget or calls (gas, fees and time still cost). A public attempt budget bounds the decision channel to $m$ bits of average min-entropy
        leakage, under assumptions we state in the theorem. On realistic cap distributions the bound
        leaves little: a vendor with five probes prices within about 4\% of the cap
        (Sections~\ref{sec:attempts}, \ref{sec:leakthm} and~\ref{sec:leakeval}).
  \item \emph{A matched comparison.} Plain x402, the same escrow path with plaintext on-chain limits,
        and the confidential design, measured interleaved on the same network, relay and vendor; an
        enclave-held policy with on-chain rollback protection, built and measured on a fork; and an
        injection study that gives every arm the same agent interface
        (Section~\ref{sec:eval}).
\end{enumerate}

\noindent We also specify a key-release protocol for digital goods that a vendor is paid for only by
publishing the decryption key, with a dispute predicate the contract decides itself; the predicate is
built and measured, and the rest of the protocol is not (Section~\ref{sec:delivery}). The live
campaign, a benchmark and a storage-wipe test ran with the signer in an Intel TDX enclave on Oasis ROFL;
the matched baselines used a signer on a local machine (Section~\ref{sec:eval}).

\paragraph{Terms.} We keep four events apart throughout. \emph{Authorization} is the vault's approval of
a request and the signer's signature over it. \emph{Funding} moves the payer's tokens into the escrow.
\emph{Capture} pays the vendor from the escrow. \emph{Delivery} is the vendor handing over what was
bought, which nothing in the implemented system ties to capture.
